% Claim proof file for row claim_3_16 -- "MargPreservesAncestors".
% Refactor twin of `claim_3_16_proof_MargPreservesAncestors.tex` produced
% under the `cdmg_typed_edges` refactor: the LN-level mathematical
% argument is unchanged, but the bidirected-edge component `L` of every
% CDMG appearing below is now encoded as a subset of `Sym2 Node` (the
% unordered-pair quotient `(V × V) / ((v1,v2) ~ (v2,v1))`) rather than
% as a `Finset (Node × Node)` paired with a separate symmetry axiom.
% Concretely this twin differs from the original only in:
%   - the statement block restatement gains a brief `\emph{Encoding of
%     bidirected edges.}` paragraph naming the `Sym2`-typing of `L`,
%     `L^{\sm W}`, and noting that swap-symmetry is now definitional and
%     irreflexivity reads as ``no self-pair'';
%   - the proof's setup paragraph names the same encoding for
%     `L^{\sm W}`, so the set-builder displayed below is read as the
%     image of the ordered-pair filter under the `Sym2` quotient map;
%   - no proof step is added, removed, or rephrased; the LN argument is
%     preserved verbatim, since the unordered-pair brace notation
%     `\lC v, w \rC \in L` (and similarly for `L^{\sm W}`) that the
%     original proof already uses for bidirected-edge membership is
%     already the quotient form and requires no further rewriting.

\documentclass[main]{subfiles}

\begin{document}
% ref: claim_3_16 (proof, with statement restated for readability)
% title: MargPreservesAncestors
\def\rowref{claim\_3\_16}
\phantomsection\label{claim_3_16}

% --- statement (restated) ---------------------------------------------------
\begin{Rem}[MargPreservesAncestors]\label{rem:marg_preserves_ancestors_bifurcations_acyclicity}
    Let $G = (J, V, E, L)$ be a CDMG (def \ref{def-cdmg}); throughout we use the notation of def \ref{not-cdmg}, the walk definitions of def \ref{def:walks}, the family-relationship notation of def \ref{def_3_5}, the acyclicity notion of def \ref{def-acylic}, the topological-order notion of def \ref{def-topological-order}, and the marginalization operator of def \ref{def:G_marginalization}. Let
    \[
        W \;\ins\; V
    \]
    be a subset of output nodes, and let
    \[
        G^{\sm W} \;=\; \lp J^{\sm W},\, V^{\sm W},\, E^{\sm W},\, L^{\sm W} \rp
    \]
    be the marginalization of $G$ with respect to $W$ in the sense of def \ref{def:G_marginalization}. Recall from items~i.\ and~ii.\ of def \ref{def:G_marginalization} that
    \[
        J^{\sm W} \;=\; J
        \qquad\text{and}\qquad
        V^{\sm W} \;=\; V \sm W;
    \]
    consequently the node set (carrier) of $G^{\sm W}$ is
    \[
        J^{\sm W} \cup V^{\sm W} \;=\; J \cup (V \sm W) \;=\; (J \cup V) \sm W,
    \]
    where the second equality uses $W \ins V$ together with $J \cap V = \emptyset$ from def \ref{def-cdmg}. Both quantifiers ``for every CDMG $G$'' and ``for every $W \ins V$'' are read \emph{universally}, with the side condition $W \ins V$ inherited verbatim from def \ref{def:G_marginalization} (so $W$ is required to consist of output nodes only); the case $W = \emptyset$ is admitted (then $V \sm W = V$, so $G^{\sm W} = G$ verbatim by items~i.--iv.\ of def \ref{def:G_marginalization}).

    \emph{Encoding of bidirected edges (refactor twin).}
    Under the project's \texttt{cdmg\_typed\_edges} refactor (the encoding in which this twin proof is written), the bidirected-edge component $L$ of every CDMG below is taken as a subset of
    \[
        \text{Sym}_2(V) \;=\; V \times V / ((v_1, v_2) \sim (v_2, v_1))
    \]
    --- the unordered-pair quotient of def \ref{def-cdmg}, realised in the project's Lean encoding by Mathlib's $\text{Sym}_2$ type. Three downstream consequences hold uniformly throughout the proof below: (i) swap-symmetry $(v_1, v_2) \sim (v_2, v_1)$ is \emph{definitional} (so the LN's separate symmetry implication $(v_1, v_2) \in L \Rightarrow (v_2, v_1) \in L$ is vacuous and not separately invoked); (ii) irreflexivity reads as ``no self-pair'' $\lC v, v \rC \notin L$ rather than as the LN's separate axiom $(v_1, v_2) \in L \Rightarrow v_1 \neq v_2$; (iii) the same encoding propagates to the marginalized component $L^{\sm W} \ins \text{Sym}_2(V \sm W)$ via the natural lift of the def \ref{def:G_marginalization} set-builder (made explicit in the proof's setup paragraph below). Concretely, the unordered-pair brace shape $\lC v_1, v_2 \rC \in L$ (resp.\ $\lC v_1, v_2 \rC \in L^{\sm W}$) for bidirected-edge membership is the canonical (and only) presentation throughout this twin; the ordered-pair shape $(v_1, v_2) \in L$ used by the pre-refactor proof is read as its image under the quotient map $V \times V \twoheadrightarrow \text{Sym}_2(V)$ whenever it appears.

    \emph{Marginalization preserves ancestral relations, bifurcations, and acyclicity:} the following three sub-claims hold.

    \begin{enumerate}[label=\roman*.)]
        \item \emph{Preservation of ancestral relations.}
            For every pair of nodes
            \[
                v_1,\, v_2 \;\in\; (J \cup V) \sm W \;=\; J \cup (V \sm W),
            \]
            the biconditional
            \[
                v_1 \in \Anc^G(v_2) \;\Longleftrightarrow\; v_1 \in \Anc^{G^{\sm W}}(v_2)
            \]
            holds, where $\Anc^G(v_2)$ is the ancestor set of $v_2$ in $G$ in the sense of def \ref{def_3_5}, item~iv., evaluated on the carrier $J \cup V$, and $\Anc^{G^{\sm W}}(v_2)$ is the ancestor set of $v_2$ in $G^{\sm W}$, also in the sense of def \ref{def_3_5}, item~iv., applied to the CDMG $G^{\sm W}$ and so evaluated on the carrier $J \cup (V \sm W)$.

            \emph{Reading of the precondition.} The LN's literal phrasing ``$v_1, v_2 \in G$ with $v_1, v_2 \notin W$'' is read set-theoretically as $v_1, v_2 \in (J \cup V) \sm W$, equivalently $v_1, v_2 \in J \cup (V \sm W)$ (the two are the same set because $W \ins V$ and $J \cap V = \emptyset$, so $J \sm W = J$). Equivalently, $v_1, v_2$ are required to lie in the carrier of $G^{\sm W}$; this is the precondition under which the right-hand membership $v_1 \in \Anc^{G^{\sm W}}(v_2)$ is well-typed.

        \item \label{rem:marg_preserves_bifurcations}
            \emph{Preservation of bifurcations (two biconditionals).}
            The following two assertions (a) and (b) both hold. They correspond to reading the LN's literal phrasing
            \begin{quote}
                ``there is a bifurcation between $v_1$ and $v_2$ (with source $v_3$) in $G$ if and only if there is a bifurcation between $v_1$ and $v_2$ (with source $v_3$) in $G^{\sm W}$''
            \end{quote}
            once with the two parenthetical clauses ``(with source $v_3$)'' omitted (this gives (a) below) and once with them included (this gives (b) below).

            \begin{enumerate}[label=(\alph*)]
                \item \emph{Sourceless biconditional.}
                    For every pair of nodes
                    \[
                        v_1,\, v_2 \;\in\; (J \cup V) \sm W \;=\; J \cup (V \sm W),
                    \]
                    the following two propositions are equivalent.

                    \emph{Reading of the precondition.} The LN's literal phrasing ``$v_1, v_2 \in G \sm W$'' is read set-theoretically as $v_1, v_2 \in (J \cup V) \sm W$, equivalently $v_1, v_2 \in J \cup (V \sm W)$, matching the outer-quantifier convention adopted in item~i.\ above and in item~iii.\ below (and matching the carrier of $G^{\sm W}$). Although the LN's $v_1, v_2$ are intended as the endpoints of the bifurcation walk and are therefore \emph{automatically} forced into $V \cap ((J \cup V) \sm W) = V \sm W$ by clauses~(b) and~(c) of def \ref{def:walks}, item~vi.\ (which require $a_0, a_{n-1}$ to be incident to the bifurcation hinge via directed edges of $G$, so that the endpoint-targets $v_0 = w_0$ and $v_n = w_n$ lie in $V$ by $E \ins (J \cup V) \times V$), no a priori restriction to $V \sm W$ is imposed on the outer quantifier: input-node endpoints simply make both sides of the biconditional vacuously false.
                    \begin{itemize}
                        \item There exists a bifurcation between $v_1$ and $v_2$ in $G$ in the sense of def \ref{def:walks}, item~vi.: i.e.\ there exist an integer $n \ge 1$, an index $k \in \{1, \dots, n\}$, and a walk $p = (w_0, a_0, w_1, a_1, \dots, w_{n-1}, a_{n-1}, w_n)$ in $G$ (def \ref{def:walks}, item~i., with $w_0, w_1, \dots, w_n \in J \cup V$ and $a_0, a_1, \dots, a_{n-1} \in (J \cup V) \times (J \cup V)$) with $\lC w_0, w_n \rC = \lC v_1, v_2 \rC$, such that $p$ together with the index $k$ satisfies clauses~(a)--(e) of def \ref{def:walks}, item~vi.
                        \item There exists a bifurcation between $v_1$ and $v_2$ in $G^{\sm W}$ in the sense of def \ref{def:walks}, item~vi., applied to the CDMG $G^{\sm W}$: i.e.\ there exist an integer $n' \ge 1$, an index $k' \in \{1, \dots, n'\}$, and a walk $p' = (w'_0, a'_0, w'_1, \dots, a'_{n' - 1}, w'_{n'})$ in $G^{\sm W}$ (with $w'_0, w'_1, \dots, w'_{n'} \in J \cup (V \sm W)$ and $a'_0, a'_1, \dots, a'_{n' - 1} \in (J \cup (V \sm W)) \times (J \cup (V \sm W))$) with $\lC w'_0, w'_{n'} \rC = \lC v_1, v_2 \rC$, such that $p'$ together with $k'$ satisfies clauses~(a)--(e) of def \ref{def:walks}, item~vi., where every occurrence of ``$E$'' (resp.\ ``$L$'') in those clauses is read as ``$E^{\sm W}$'' (resp.\ ``$L^{\sm W}$'').
                    \end{itemize}
                    The set equality $\lC w_0, w_n \rC = \lC v_1, v_2 \rC$ (and analogously $\lC w'_0, w'_{n'} \rC = \lC v_1, v_2 \rC$) makes the symmetric ``between $v_1$ and $v_2$'' phrasing explicit: either $w_0 = v_1, w_n = v_2$ or $w_0 = v_2, w_n = v_1$ is admissible.

                \item \emph{Sourced biconditional.}
                    For every triple of nodes
                    \[
                        v_1,\, v_2,\, v_3 \;\in\; (J \cup V) \sm W \;=\; J \cup (V \sm W),
                    \]
                    the following two propositions are equivalent.

                    \emph{Reading of the precondition.} The LN's literal phrasing ``$v_1, v_2, v_3 \in G \sm W$'' is again read set-theoretically as $v_1, v_2, v_3 \in (J \cup V) \sm W = J \cup (V \sm W)$, matching items~i.\ and~iii.\ As in sub-claim~(a), the bifurcation endpoints $v_1, v_2$ are automatically forced into $V \sm W$ by clauses~(b) and~(c) of def \ref{def:walks}, item~vi.; the source $v_3 = w_k$, however, is \emph{not} so forced --- by clause~(d) directed alternative $a_{k-1} = (w_k, w_{k-1}) \in E$ together with $E \ins (J \cup V) \times V$, only $w_{k-1} \in V$ is implied, while $w_k$ may lie in $J$ (cf.\ prp \ref{prp:bifurcations_alternative}, which types the source $c \in J \cup V$). Hence the source $v_3$ may genuinely be an input node $v_3 \in J$, and the broader outer quantifier $(J \cup V) \sm W$ on $v_3$ is non-vacuously needed.
                    \begin{itemize}
                        \item There exists a bifurcation between $v_1$ and $v_2$ with source $v_3$ in $G$, in the sense of def \ref{def:walks}, item~vi., and its closing paragraph ``Source of the bifurcation'': i.e.\ there exist an integer $n \ge 1$, an index $k \in \{1, \dots, n\}$, and a walk $p = (w_0, a_0, w_1, \dots, a_{n-1}, w_n)$ in $G$ with $\lC w_0, w_n \rC = \lC v_1, v_2 \rC$, satisfying clauses~(a)--(e) of def \ref{def:walks}, item~vi., \emph{and} clause~(d) of def \ref{def:walks}, item~vi., is realised at $k$ by the \emph{directed} alternative $a_{k-1} = (w_k, w_{k-1}) \in E$ with $w_k = v_3$ (so $v_3 = w_k$ is the source of $p$ in the sense of def \ref{def:walks}, item~vi., closing paragraph). By clause~(e) of def \ref{def:walks}, item~vi., the source $v_3 = w_k$ is automatically an interior vertex of $p$ with $1 \le k \le n - 1$ (hence $n \ge 2$), and $v_3 \ne v_1$, $v_3 \ne v_2$.
                        \item There exists a bifurcation between $v_1$ and $v_2$ with source $v_3$ in $G^{\sm W}$, in the sense of def \ref{def:walks}, item~vi., applied to the CDMG $G^{\sm W}$: i.e.\ the analogous existential, with primed objects $n', k', p' = (w'_0, a'_0, \dots, a'_{n' - 1}, w'_{n'})$ in $G^{\sm W}$, all walk vertices $w'_0, \dots, w'_{n'}$ drawn from the carrier $J \cup (V \sm W)$ of $G^{\sm W}$, $\lC w'_0, w'_{n'} \rC = \lC v_1, v_2 \rC$, clauses~(a)--(e) of def \ref{def:walks}, item~vi., satisfied with ``$E$'' read as ``$E^{\sm W}$'' and ``$L$'' read as ``$L^{\sm W}$'', and clause~(d) realised at $k'$ by the directed alternative $a'_{k' - 1} = (w'_{k'}, w'_{k' - 1}) \in E^{\sm W}$ with $w'_{k'} = v_3$.
                    \end{itemize}
            \end{enumerate}
            The biconditional in (a) is the LN's ``parenthetical clauses omitted'' reading; the biconditional in (b) is the LN's ``parenthetical clauses included'' reading. The addition to the LN explicitly authorises asserting both.

        \item \emph{Preservation of acyclicity and induced topological order.}
            Suppose $G$ is acyclic in the sense of def \ref{def-acylic}. Then both of the following sub-assertions (a) and (b) hold.
            \begin{enumerate}[label=(\alph*)]
                \item \emph{Acyclicity is preserved.}
                    $G^{\sm W}$ is acyclic in the sense of def \ref{def-acylic}, evaluated on the carrier $J \cup (V \sm W)$ of $G^{\sm W}$: explicitly, for every $v \in J \cup (V \sm W)$, there does not exist any non-trivial directed walk from $v$ to $v$ in $G^{\sm W}$ (where ``directed walk in $G^{\sm W}$'' is in the sense of def \ref{def:walks}, item~ii., applied to $G^{\sm W}$ --- with edge-membership read against $E^{\sm W}$ instead of $E$ --- and ``non-trivial'' means walk length $n \ge 1$).

                \item \emph{Topological order is induced by restriction.}
                    For every binary relation $<_G$ on $J \cup V$ that is a topological order of $G$ in the sense of def \ref{def-topological-order}, the binary relation $<_{G^{\sm W}}$ on $J \cup (V \sm W)$ defined by \emph{restriction of $<_G$ to the smaller carrier}, namely
                    \[
                        \forall\, v_1, v_2 \in J \cup (V \sm W) :\quad
                        v_1 \,<_{G^{\sm W}}\, v_2 \;\;:\Longleftrightarrow\;\; v_1 \,<_G\, v_2,
                    \]
                    is a topological order of $G^{\sm W}$ in the sense of def \ref{def-topological-order}, evaluated on the carrier $J \cup (V \sm W)$ of $G^{\sm W}$.

                    \emph{Reading of the LN's ``induces'' / ``by just ignoring the nodes from $W$''.}
                    The LN's clause ``a topological order of $G$ induces a topological order on $G^{\sm W}$ (by just ignoring the nodes from $W$)'' is read here as the explicit \emph{restriction operation} above: the carrier of $<_{G^{\sm W}}$ is the smaller set $J \cup (V \sm W) = (J \cup V) \sm W$, the input nodes $J$ are unchanged (since $W \ins V$ does not intersect $J$ by $J \cap V = \emptyset$), the output nodes $W$ are removed from the carrier, and for every pair $(v_1, v_2)$ both of whose entries lie in the smaller carrier, the value of $<_{G^{\sm W}}$ on $(v_1, v_2)$ coincides with the value of $<_G$ on the same pair. The LN's ``by just ignoring the nodes from $W$'' is the assertion that the value of $<_{G^{\sm W}}(v_1, v_2)$ does not require any reassignment of order between unsplit nodes; it equals $<_G(v_1, v_2)$ on every pair in the smaller carrier.

                    Being a topological order of $G^{\sm W}$ means, per def \ref{def-topological-order} applied to $G^{\sm W}$, that $<_{G^{\sm W}}$ is a strict total order on $J \cup (V \sm W)$ (i.e.\ irreflexive, transitive, and trichotomous on $J \cup (V \sm W)$) and that for every $v_1, v_2 \in J \cup (V \sm W)$, $v_1 \in \Pa^{G^{\sm W}}(v_2)$ implies $v_1 <_{G^{\sm W}} v_2$, where $\Pa^{G^{\sm W}}(v_2) = \lC u \in J \cup (V \sm W) \st (u, v_2) \in E^{\sm W} \rC$ is the parent set of $v_2$ in $G^{\sm W}$ in the sense of def \ref{def_3_5}, item~i.
            \end{enumerate}
    \end{enumerate}
\end{Rem}

% --- proof ------------------------------------------------------------------
\begin{proof}
% The LN block in lecture-notes/lecture_notes/graphs.tex:962--975 is a
% \begin{Rem} (three numbered sub-items) followed by a \begin{proof} at
% lines 976--993 that addresses ONLY sub-claim ii of the LN's literal
% three-item list (the bifurcation biconditional, in its single-claim
% reading). Sub-claims i and iii of the LN's literal list have no LN-supplied
% proof at all, and the LN's proof of sub-claim ii uses induction on |W|
% together with \refrow{claim_3_17} (marginalizations commute). Because
% \refrow{claim_3_17} sits AFTER \refrow{claim_3_16} in the chapter row order
% (it is a downstream consumer of the present row), we cannot import that
% induction; we replace it with the LN's converse-direction
% projection / expansion argument, applied uniformly in both directions and
% to all five sub-assertions of the rewritten statement above.

We use throughout the expanded notation of the statement block above. We open with a unifying setup paragraph that names the two witness predicates of items~iii.\ and~iv.\ of def \ref{def:G_marginalization} and pins down how they are used in projection / expansion arguments; the five sub-claims then follow in order.

\medskip\noindent\emph{Setup: the marginalization witness predicates.}
Per items~iii.\ and~iv.\ of def \ref{def:G_marginalization},
\begin{align*}
    E^{\sm W} \;&=\; \lC (\ul{v}, \ol{v}) \in (J \cup (V \sm W)) \times (V \sm W) \st \Phi_E^{\sm W}(\ul{v}, \ol{v}) \rC,\\
    L^{\sm W} \;&=\; \lC \lC \ul{v}, \ol{v} \rC \in \text{Sym}_2(V \sm W) \st \ul{v} \neq \ol{v} \;\wedge\; \Phi_L^{\sm W}(\ul{v}, \ol{v}) \rC,
\end{align*}
where the bidirected component $L^{\sm W}$ is read under the \texttt{cdmg\_typed\_edges} encoding announced in the statement block (so $L^{\sm W} \ins \text{Sym}_2(V \sm W)$ is the natural lift of the def \ref{def:G_marginalization} ordered-pair filter under the quotient map $(V \sm W) \times (V \sm W) \twoheadrightarrow \text{Sym}_2(V \sm W)$; the conjunct $\Phi_L^{\sm W}(\ul{v}, \ol{v})$ is well-defined on the unordered pair because, as $\Phi_L^{\sm W}$ asserts the existence of a bifurcation between $\ul{v}$ and $\ol{v}$ in $G$ --- and bifurcations are symmetric in their two end-vertices by def \ref{def:walks}, item~vi.\ --- it descends through the quotient verbatim), and the membership predicates $\Phi_E^{\sm W}$ and $\Phi_L^{\sm W}$ unfold as follows:
\begin{itemize}
    \item $\Phi_E^{\sm W}(\ul{v}, \ol{v})$ asserts the existence of a directed walk in $G$ (in the sense of def \ref{def:walks}, item~ii.) from $\ul{v}$ to $\ol{v}$ of length $n \ge 1$ whose intermediate vertices $w_1, \dots, w_{n-1}$ (if any) all lie in $W$. (Note: per the LN footnote at item~iii.\ of def \ref{def:G_marginalization}, self-cycles $\ul{v} = \ol{v}$ are admitted, in particular by length-$1$ direct edges $(\ul{v}, \ul{v}) \in E$ with $\ul{v} \in V \sm W$, and equally by longer $W$-traversing walks $\ul{v} \to w_1 \to \dots \to \ul{v}$.)
    \item $\Phi_L^{\sm W}(\ul{v}, \ol{v})$ asserts the existence of a bifurcation in $G$ (in the sense of def \ref{def:walks}, item~vi.) between $\ul{v}$ and $\ol{v}$ --- i.e.\ a walk $\tilde{p}$ together with a split index $\tilde{k}$ satisfying clauses~(a)--(e) --- with all intermediate vertices of $\tilde{p}$ (the vertices $w_1, \dots, w_{n-1}$ in $\tilde{p}$'s expansion) lying in $W$.
\end{itemize}
The witness predicates are forward-only / backward-only existentials: writing them down in either direction is just an unfolding of the set-builder definitions in items~iii.\ and~iv.\ of def \ref{def:G_marginalization}. We will use the two directions repeatedly below: the \emph{contraction direction} (a witnessing $G$-walk witnesses membership of a pair in $E^{\sm W}$ or $L^{\sm W}$) and the \emph{expansion direction} (membership of a pair in $E^{\sm W}$ or $L^{\sm W}$ gives back a witnessing $G$-walk).

We address the five sub-assertions of the statement block (i, ii(a), ii(b), iii(a), iii(b)) in turn. The projection / expansion duality applied in sub-claim~i extends, with the same structural moves but applied to bifurcation walks instead of directed walks, to sub-claims~ii(a) and~ii(b); sub-claims~iii(a) and~iii(b) are then short corollaries of sub-claim~i's expansion direction.

\medskip\noindent\textbf{Sub-claim~i: preservation of ancestral relations.}

Let $v_1, v_2 \in (J \cup V) \sm W = J \cup (V \sm W)$. By def \ref{def_3_5}, item~iv., the ancestor predicate $v_1 \in \Anc^G(v_2)$ unfolds (on the carrier $J \cup V$) to the conjunction of $v_1 \in J \cup V$ and ``there exists a directed walk in $G$ from $v_1$ to $v_2$'' in the sense of def \ref{def:walks}, item~ii.\ (with length-zero walks $(v_1)$ admitted, witnessing self-ancestry $v_1 \in \Anc^G(v_1)$). The analogous unfolding for $v_1 \in \Anc^{G^{\sm W}}(v_2)$ yields the conjunction of $v_1 \in J \cup (V \sm W)$ and ``there exists a directed walk in $G^{\sm W}$ from $v_1$ to $v_2$'' (edge-membership read against $E^{\sm W}$). The membership preconditions $v_1 \in J \cup V$ and $v_1 \in J \cup (V \sm W)$ are both granted by hypothesis, since $v_1 \in (J \cup V) \sm W = J \cup (V \sm W) \subseteq J \cup V$. Sub-claim~i therefore reduces to the equivalence
\[
    \exists \text{ directed walk in } G \text{ from } v_1 \text{ to } v_2
    \;\iff\;
    \exists \text{ directed walk in } G^{\sm W} \text{ from } v_1 \text{ to } v_2.
\]

\smallskip\noindent\emph{($\Rightarrow$) Forward direction (contraction).}
Let $p = (u_0, a_0, u_1, a_1, \dots, a_{m-1}, u_m)$ be a directed walk in $G$ with $u_0 = v_1$, $u_m = v_2$. In particular $a_i = (u_i, u_{i+1}) \in E$ for every $i \in \{0, 1, \dots, m-1\}$ (def \ref{def:walks}, item~ii.). We construct a directed walk in $G^{\sm W}$ from $v_1$ to $v_2$ by projecting $p$ onto its non-$W$ vertices.

\emph{Degenerate sub-case $m = 0$ (trivial walk).} Then $v_1 = u_0 = u_m = v_2$, and the trivial walk $(v_1)$ in $G^{\sm W}$ (well-typed because $v_1 \in J \cup (V \sm W)$, the carrier of $G^{\sm W}$) is a directed walk in $G^{\sm W}$ from $v_1$ to $v_2 = v_1$ of length $0$, witnessing the RHS.

\emph{Main sub-case $m \ge 1$.}
Enumerate, in strictly increasing order, the indices $i \in \{0, 1, \dots, m\}$ for which $u_i \notin W$:
\[
    0 \;=\; c_0 \;<\; c_1 \;<\; \cdots \;<\; c_s \;=\; m.
\]
The endpoint indices $0$ and $m$ are included because $u_0 = v_1 \notin W$ and $u_m = v_2 \notin W$ (by hypothesis $v_1, v_2 \in (J \cup V) \sm W$). For each $j \in \{0, 1, \dots, s-1\}$, consider the sub-walk of $p$ from index $c_j$ to index $c_{j+1}$:
\[
    \tilde{p}_j \;:=\; \lp u_{c_j},\, a_{c_j},\, u_{c_j + 1},\, a_{c_j + 1},\, \dots,\, a_{c_{j+1} - 1},\, u_{c_{j+1}} \rp.
\]
This $\tilde{p}_j$ is a directed walk in $G$ from $u_{c_j}$ to $u_{c_{j+1}}$ of length $c_{j+1} - c_j \ge 1$, and all its intermediate vertices $u_{c_j + 1}, u_{c_j + 2}, \dots, u_{c_{j+1} - 1}$ lie in $W$ by maximality of the non-$W$ subsequence (any non-$W$ index between $c_j$ and $c_{j+1}$ would have been included as some $c_{j'}$).

Pair typing: $u_{c_j} \in J \cup V$ (as a walk vertex of $G$, per def \ref{def-cdmg}'s carrier) and $u_{c_j} \notin W$, so $u_{c_j} \in J \cup (V \sm W)$; for the head, $u_{c_{j+1}}$ is the target of the directed edge $a_{c_{j+1} - 1} \in E$, so $u_{c_{j+1}} \in V$ by $E \ins (J \cup V) \times V$ (def \ref{def-cdmg}), and $u_{c_{j+1}} \notin W$ gives $u_{c_{j+1}} \in V \sm W$. Hence the pair $(u_{c_j}, u_{c_{j+1}})$ lies in the product carrier $(J \cup (V \sm W)) \times (V \sm W)$ of $E^{\sm W}$.

Apply $\Phi_E^{\sm W}$ to $\tilde{p}_j$: the walk $\tilde{p}_j$ witnesses the existential of $\Phi_E^{\sm W}(u_{c_j}, u_{c_{j+1}})$ directly --- length $\ge 1$ ✓, intermediates in $W$ ✓, the length-$\ge 1$ requirement is met directly by $c_{j+1} - c_j \ge 1$. By the contraction direction of $\Phi_E^{\sm W}$ (item~iii.\ of def \ref{def:G_marginalization}), $(u_{c_j}, u_{c_{j+1}}) \in E^{\sm W}$.

Concatenate the projected edges $(u_{c_j}, u_{c_{j+1}}) \in E^{\sm W}$ for $j = 0, 1, \dots, s-1$ along their shared endpoints $u_{c_1}, u_{c_2}, \dots, u_{c_{s-1}}$ (after omitting the zero-net-step degenerate self-loops mentioned above) to obtain a directed walk in $G^{\sm W}$ from $u_{c_0} = v_1$ to $u_{c_s} = v_2$ in the sense of def \ref{def:walks}, item~ii., applied to $G^{\sm W}$.

\smallskip\noindent\emph{($\Leftarrow$) Converse direction (expansion).}
Let $p' = (u'_0, a'_0, u'_1, a'_1, \dots, a'_{s-1}, u'_s)$ be a directed walk in $G^{\sm W}$ with $u'_0 = v_1$, $u'_s = v_2$, so each $a'_j = (u'_j, u'_{j+1}) \in E^{\sm W}$ (def \ref{def:walks}, item~ii., applied to $G^{\sm W}$).

\emph{Degenerate sub-case $s = 0$.} Then $v_1 = u'_0 = u'_s = v_2$, and the trivial walk $(v_1)$ in $G$ (well-typed because $v_1 \in J \cup V$, the carrier of $G$, by $v_1 \in (J \cup V) \sm W \subseteq J \cup V$) is a directed walk in $G$ of length $0$ from $v_1$ to $v_2 = v_1$, witnessing the LHS.

\emph{Main sub-case $s \ge 1$.} For each $j \in \{0, 1, \dots, s-1\}$, $(u'_j, u'_{j+1}) \in E^{\sm W}$, so by the expansion direction of $\Phi_E^{\sm W}$ (item~iii.\ of def \ref{def:G_marginalization}) there exists a directed walk $q_j$ in $G$ from $u'_j$ to $u'_{j+1}$ of length $\ge 1$ with all intermediate vertices in $W$. Concatenate $q_0, q_1, \dots, q_{s-1}$ along their shared endpoints $u'_1, u'_2, \dots, u'_{s-1}$ to obtain a directed walk in $G$ from $u'_0 = v_1$ to $u'_s = v_2$ in the sense of def \ref{def:walks}, item~ii.\ This is the witnessing $G$-walk for the LHS.

This establishes sub-claim~i.

\medskip\noindent\textbf{Sub-claim~ii(a): sourceless bifurcation biconditional.}

Let $v_1, v_2 \in (J \cup V) \sm W = J \cup (V \sm W)$. We show the equivalence of the two existentials of sub-claim~ii(a) by the same projection / expansion duality used for sub-claim~i, but with bifurcation walks in place of directed walks.

\smallskip\noindent\emph{($\Rightarrow$) Forward direction (contraction of a bifurcation in $G$).}
Let $p = (w_0, a_0, w_1, a_1, \dots, a_{n-1}, w_n)$ together with a split index $k \in \{1, \dots, n\}$ and length $n \ge 1$ witness a bifurcation between $v_1$ and $v_2$ in $G$ per def \ref{def:walks}, item~vi.: clauses~(a)--(e) hold for $p$ and $k$, and $\lC w_0, w_n \rC = \lC v_1, v_2 \rC$. Without loss of generality $w_0 = v_1$, $w_n = v_2$ (the other orientation $w_0 = v_2$, $w_n = v_1$ produces the conclusion under the symmetric disjunct of the RHS existential).

Recall from def \ref{def:walks}, item~vi., that the edges of $p$ are typed as follows (clauses~(b)--(d)):
\begin{itemize}
    \item Left-arm edges: $a_i = (w_{i+1}, w_i) \in E$ for $0 \le i \le k-2$ (i.e.\ directed from $w_{i+1}$ toward $w_0$).
    \item Hinge: $a_{k-1}$ is either bidirected $\lC w_{k-1}, w_k \rC \in L$, or directed $a_{k-1} = (w_k, w_{k-1}) \in E$ (the latter is the case with source $w_k$).
    \item Right-arm edges: $a_i = (w_i, w_{i+1}) \in E$ for $k \le i \le n-1$ (i.e.\ directed from $w_i$ toward $w_n$).
\end{itemize}
By clauses~(b) and~(c) of def \ref{def:walks}, item~vi., together with $E \ins (J \cup V) \times V$ from def \ref{def-cdmg} applied at the boundary edges $a_0$ and $a_{n-1}$, the endpoints $w_0 = v_1$ and $w_n = v_2$ lie in $V$. By hypothesis they also lie outside $W$, so $w_0, w_n \in V \sm W \subseteq J \cup (V \sm W)$.

\emph{The projection.} Enumerate, in strictly increasing order, the indices $i \in \{0, 1, \dots, n\}$ for which $w_i \notin W$:
\[
    0 \;=\; c_0 \;<\; c_1 \;<\; \cdots \;<\; c_s \;=\; n.
\]
(The endpoint indices $0$ and $n$ are included because $w_0, w_n \notin W$, hence $s \ge 0$.) Set $w'_j := w_{c_j}$ for $0 \le j \le s$, so $w'_0 = v_1$ and $w'_s = v_2$. The projected sequence $(w'_0, w'_1, \dots, w'_s)$ has length $s$ as a walk.

\emph{The projected split index.} There exists a unique index $k' \in \{1, \dots, s\}$ such that $c_{k'-1} \le k-1$ and $c_{k'} \ge k$: existence holds because $c_0 = 0 \le k-1$ (so the set $\{j : c_j \le k-1\}$ is non-empty) and $c_s = n \ge k$ (so $k' \le s$); uniqueness holds because $(c_j)$ is strictly increasing, so there is exactly one ``crossing index''. Geometrically, $w'_{k'-1} = w_{c_{k'-1}}$ is the last non-$W$ vertex on the left arm or at the hinge-left position $w_{k-1}$, and $w'_{k'} = w_{c_{k'}}$ is the first non-$W$ vertex at the hinge-right position $w_k$ or on the right arm.

\emph{Constructing the projected edges $a'_0, \dots, a'_{s-1}$.}
For each $j \in \{0, 1, \dots, s-1\}$, we choose $a'_j$ between $w'_j$ and $w'_{j+1}$ from $E^{\sm W} \cup L^{\sm W}$ according to which of three regions the consecutive pair $(c_j, c_{j+1})$ lies in:

\emph{Region L --- left arm, $c_{j+1} \le k-1$ (so $j \le k' - 2$).}
The sub-walk of $p$ from index $c_j$ to index $c_{j+1}$ consists exclusively of left-arm edges $a_{c_j}, a_{c_j + 1}, \dots, a_{c_{j+1} - 1}$, each of the form $a_i = (w_{i+1}, w_i) \in E$. Read in the reverse direction (from index $c_{j+1}$ down to $c_j$), these edges form a directed walk
\[
    \lp w_{c_{j+1}},\, a_{c_{j+1} - 1},\, w_{c_{j+1} - 1},\, \dots,\, a_{c_j},\, w_{c_j} \rp
\]
in $G$ from $w_{c_{j+1}}$ to $w_{c_j}$ of length $c_{j+1} - c_j \ge 1$, with all intermediate vertices in $W$ (by maximality of the non-$W$ subsequence). By the contraction direction of $\Phi_E^{\sm W}$ applied to this sub-walk, $(w_{c_{j+1}}, w_{c_j}) \in E^{\sm W}$. Take
\[
    a'_j \;:=\; (w'_{j+1}, w'_j) \;\in\; E^{\sm W},
\]
a left-arm edge of $p'$ (directed from $w'_{j+1}$ toward $w'_j$).

\emph{Region R --- right arm, $c_j \ge k$ (so $j \ge k'$).}
The sub-walk of $p$ from index $c_j$ to index $c_{j+1}$ consists exclusively of right-arm edges $a_{c_j}, a_{c_j + 1}, \dots, a_{c_{j+1} - 1}$, each of the form $a_i = (w_i, w_{i+1}) \in E$. Concatenated, these form a directed walk
\[
    \lp w_{c_j},\, a_{c_j},\, w_{c_j + 1},\, \dots,\, a_{c_{j+1} - 1},\, w_{c_{j+1}} \rp
\]
in $G$ from $w_{c_j}$ to $w_{c_{j+1}}$ of length $\ge 1$ with all intermediate vertices in $W$. By the contraction direction of $\Phi_E^{\sm W}$, $(w_{c_j}, w_{c_{j+1}}) \in E^{\sm W}$. Take
\[
    a'_j \;:=\; (w'_j, w'_{j+1}) \;\in\; E^{\sm W},
\]
a right-arm edge of $p'$.

\emph{Region H --- hinge straddle, $c_j \le k - 1$ and $c_{j+1} \ge k$ (so $j = k' - 1$).}
The sub-walk of $p$ from index $c_j$ to index $c_{j+1}$ contains the hinge edge $a_{k-1}$ together with portions of both arms. Its overall shape:
\[
    \lp w_{c_j},\, a_{c_j},\, w_{c_j + 1},\, \dots,\, a_{k-2},\, w_{k-1},\, a_{k-1},\, w_k,\, a_k,\, \dots,\, a_{c_{j+1} - 1},\, w_{c_{j+1}} \rp,
\]
where the prefix $(w_{c_j}, a_{c_j}, \dots, w_{k-1})$ is a portion of $p$'s left arm (reverse-directed: edges $a_i = (w_{i+1}, w_i) \in E$ for $c_j \le i \le k-2$), the middle $a_{k-1}$ is $p$'s hinge (of the same type as $p$ --- bidirected $\lC w_{k-1}, w_k \rC \in L$ or directed $(w_k, w_{k-1}) \in E$), and the suffix $(w_k, a_k, \dots, w_{c_{j+1}})$ is a portion of $p$'s right arm (directed: edges $a_i = (w_i, w_{i+1}) \in E$ for $k \le i \le c_{j+1} - 1$). All intermediate vertices of this hinge-straddling sub-walk --- namely $w_{c_j + 1}, \dots, w_{k-1}, w_k, \dots, w_{c_{j+1} - 1}$ minus the hinge vertices themselves and the endpoints --- lie in $W$ by maximality of the non-$W$ subsequence.

% correction (2026-06-12): The construction below was rewritten to repair
% two gaps surfaced by verify_tex_proof. (1) The original argument inherited
% $p$'s bifurcation structure into the hinge-straddle sub-walk ``at the same
% hinge position''. This is unsound when $p$'s hinge is directed with source
% $w_k \notin W$: then $c_{j+1} = k$, the sub-walk's local hinge index
% $\tilde k = k - c_j$ equals the sub-walk's length
% $\tilde n = c_{j+1} - c_j$, and clause~(e) of def \ref{def:walks},
% item~vi., \emph{excludes} the directed alternative at $\tilde k = \tilde n$
% (the boundary configuration ``$k = n$ with directed hinge'' is exactly
% what clause~(e) was added to rule out). A case-split on $p$'s hinge type
% plus a directed-hinge $p'$-construction (paralleling sub-claim~ii(b)'s
% Region~H) repairs this. (2) The degenerate identification
% $w_{c_j} = w_{c_{j+1}}$, previously handled by deleting a duplicate
% projection index --- unsound because the deleted vertex sits as a
% non-$W$ vertex in the interior of the resulting projection sub-walk,
% violating the intermediates-in-$W$ requirement of $\Phi_E^{\sm W}$ /
% $\Phi_L^{\sm W}$ --- is repaired by pre-processing $p$ into a shorter
% bifurcation $p^*$ in $G$ via cycle excision; the refreshed projection
% then falls under sub-case~(H.B) below.

We handle Region~H by a pre-processing step for the corner case
$w_{c_j} = w_{c_{j+1}}$, followed by a case-split (assuming
$w_{c_j} \ne w_{c_{j+1}}$) on the type of $p$'s hinge $a_{k - 1}$ together
with the placement of $p$'s hinge vertex $w_k$ with respect to $W$.

\emph{Pre-processing the degenerate identification $w_{c_j} = w_{c_{j+1}}$.}
Suppose $w_{c_j} = w_{c_{j+1}} =: v^*$. Clauses~(b) and (c) for $p$
($w_0 = v_1$ appears at index $0$ only, $w_n = v_2$ at index $n$ only)
force $1 \le c_j$ and $c_{j+1} \le n - 1$ (both strictly interior of
$p$), so $v^* \notin \{v_1, v_2\}$. The hinge-straddle sub-walk of $p$
from $c_j$ to $c_{j+1}$ then traces a closed loop $v^* \to \dots \to v^*$
whose $\tilde n - 1 = c_{j+1} - c_j - 1$ interior vertices
$w_{c_j + 1}, \dots, w_{c_{j+1} - 1}$ all lie in $W$ by maximality of
the non-$W$ subsequence (and the loop contains $p$'s hinge edge
$a_{k - 1}$ in its interior). Define a shorter walk $p^*$ in $G$ by
\emph{excising this loop} from $p$:
\[
    p^* \;:=\; \lp w_0,\, a_0,\, \dots,\, a_{c_j - 1},\, w_{c_j} = v^*,\, a_{c_{j+1}},\, w_{c_{j+1} + 1},\, \dots,\, a_{n - 1},\, w_n \rp,
\]
of length $n^* := n - (c_{j+1} - c_j)$. Place the hinge of $p^*$ at
$p^*$-index $k^* := c_j$ (the position of $v^*$ in $p^*$'s vertex list).
At this position the two incident walk-edges are
$a_{c_j - 1} = (w_{c_j}, w_{c_j - 1}) = (v^*, w_{c_j - 1}) \in E$ (from
$p$'s left arm, since $c_j - 1 \le k - 2$) and
$a_{c_{j+1}} = (w_{c_{j+1}}, w_{c_{j+1} + 1}) = (v^*, w_{c_{j+1} + 1}) \in E$
(from $p$'s right arm, since $c_{j+1} \ge k$), both directed \emph{out
of} $v^*$. Hence $p^*$ has a \emph{directed-fork hinge at $v^*$ with
source $v^*$}: the hinge edge $a_{k^* - 1} = a_{c_j - 1}$ realises
clause~(d)'s directed alternative
$(v_{k^*}^{p^*}, v_{k^* - 1}^{p^*}) \in E$ with $v_{k^*}^{p^*} = v^*$.

Clauses~(a)--(e) of def \ref{def:walks}, item~vi., for $p^*$ at $k^*$:
\begin{itemize}
    \item Clause~(a) is inherited (no new repetitions are introduced; $v_1, v_2$ remain at their unique endpoint positions $0$ and $n^*$ of $p^*$'s vertex list, since the excised loop did not visit $v_1$ or $v_2$ by clauses~(b), (c) for $p$).
    \item Clauses~(b), (c) are inherited from the unchanged portions of $p$: $p^*$'s left arm (indices $0, \dots, k^* - 1$, edges $a_0, \dots, a_{c_j - 2}$) keeps $p$'s left-arm orientation $a_i = (w_{i+1}, w_i) \in E$; $p^*$'s right arm (indices $k^*, \dots, n^* - 1$ in edges, $k^*, \dots, n^*$ in vertices, with edges $a_{c_{j+1}}, \dots, a_{n - 1}$) keeps $p$'s right-arm orientation $a_i = (w_i, w_{i+1}) \in E$. (The edge $a_{c_j - 1}$, formerly on $p$'s left arm, is reassigned to $p^*$'s hinge.)
    \item Clause~(d) is satisfied by the directed-fork hinge at $k^*$ above.
    \item Clause~(e) is satisfied because $k^* = c_j \ge 1$ and $n^* - k^* = n - c_{j+1} \ge 1$ place $k^*$ strictly interior to $p^*$; the boundary edges $a_0$ (still on $p^*$'s left arm) and $a_{n - 1}$ (still on $p^*$'s right arm) are unchanged and inherit $p$'s arrowhead orientations at $v_1$ and $v_2$ respectively.
\end{itemize}
So $p^*$ together with $k^*$ is a valid bifurcation in $G$ between $v_1$
and $v_2$.

Replace $p$ by $p^*$, refresh the projection enumeration on $p^*$'s
vertex list (the original $(c_i)_{i = 0}^{s}$ with $c_{j+1}$ removed and
the indices to its right shifted down by $c_{j+1} - c_j$, of length $s$,
hence $s - 1$ projection edges), and recompute the crossing index of
$p^*$'s hinge in the refreshed enumeration. Since $v^*$ sits at
$p^*$-index $c_j$ which remains a non-$W$ index of $p^*$, the refreshed
enumeration has $j$-th entry $c_j$ and $(j - 1)$-th entry $c_{j - 1}$
(with $c_{j - 1} \le c_j - 1$), so the refreshed crossing index is $j$.
The refreshed Region~H hinge-straddle is the sub-walk of $p^*$ from
refreshed projection index $c_{j - 1}$ to $c_j$, containing $p^*$'s hinge
edge $a_{c_j - 1}$. Since $p^*$'s hinge is directed with source
$v^* \notin W$, this refreshed Region~H falls under sub-case~(H.B) below
applied to $p^*$, producing a \emph{directed} $p'$-hinge edge in
$E^{\sm W}$ for the refreshed projection. (If the refreshed Region~H of
$p^*$'s projection itself exhibits a degenerate identification
$w_{c_{j - 1}} = v^*$, iterate the pre-processing on $p^*$ in place of
$p$; each iteration strictly shortens the walk by at least one edge, so
the iteration terminates with a walk whose final Region~H endpoints are
distinct and whose hinge remains directed with source $v^* \notin W$,
still in sub-case~(H.B).)

\emph{Henceforth assume $w_{c_j} \ne w_{c_{j+1}}$ in Region~H} (the
degenerate case having been pre-processed as above into sub-case~(H.B)).
We case-split on the type of $p$'s hinge.

\emph{Sub-case (H.A): $p$'s hinge is bidirected, \textbf{or} $p$'s hinge
is directed with source $w_k \in W$.}
The hinge-straddle sub-walk of $p$ from index $c_j$ to index $c_{j+1}$,
together with the local split index $\tilde k := k - c_j$ and length
$\tilde n := c_{j+1} - c_j$, is a bifurcation in $G$ between $w_{c_j}$
and $w_{c_{j+1}}$:
\begin{itemize}
    \item Clause~(a). $w_{c_j} \ne w_{c_{j+1}}$ by the henceforth assumption; the end-nodes-appearing-once condition for the sub-walk is inherited from $p$'s structure at indices $c_j, c_{j+1}$ (the sub-walk's interior vertices $w_{c_j + 1}, \dots, w_{c_{j+1} - 1}$ all lie in $W$, while $w_{c_j}, w_{c_{j+1}}$ do not, so no interior position equals either endpoint).
    \item Clauses~(b), (c). The sub-walk's left-arm edges $a_{c_j}, \dots, a_{k - 2}$ inherit $p$'s left-arm orientation $a_i = (w_{i + 1}, w_i) \in E$ (vacuous when $\tilde k = 1$); the sub-walk's right-arm edges $a_k, \dots, a_{c_{j+1} - 1}$ inherit $p$'s right-arm orientation $a_i = (w_i, w_{i + 1}) \in E$ (vacuous when $\tilde k = \tilde n$).
    \item Clause~(d). The sub-walk's hinge edge $a_{\tilde k - 1} = a_{k - 1}$ inherits $p$'s hinge type (bidirected or directed) verbatim.
    \item Clause~(e). \emph{If $p$'s hinge is bidirected}: the bidirected alternative is admitted by clause~(e) of def \ref{def:walks}, item~vi., at every $\tilde k \in \{1, \dots, \tilde n\}$ --- including the boundary $\tilde k = \tilde n$, which occurs when $w_k \notin W$ (so $c_{j+1} = k$) --- because clause~(e) only \emph{excludes} the directed alternative at $\tilde k = \tilde n$, not the bidirected one. \emph{If $p$'s hinge is directed with $w_k \in W$}: then $w_k$ does not appear in the projection enumeration, so $c_{j+1} > k$, giving $\tilde k = k - c_j < c_{j+1} - c_j = \tilde n$; combined with $\tilde k \ge 1$ (from $k \ge 1$ and $c_j \le k - 1$), the sub-walk's directed hinge sits at the strictly interior $\tilde k \in \{1, \dots, \tilde n - 1\}$, where clause~(e) admits the directed alternative.
\end{itemize}
By the contraction direction of $\Phi_L^{\sm W}$ (item~iv.\ of
def \ref{def:G_marginalization}), $\Phi_L^{\sm W}(w_{c_j}, w_{c_{j+1}})$
holds (the hinge-straddle sub-walk is itself a witnessing bifurcation
in $G$ between $w_{c_j}$ and $w_{c_{j+1}}$ with intermediate vertices in
$W$). Combined with $w_{c_j} \ne w_{c_{j+1}}$,
$\lC w_{c_j}, w_{c_{j+1}} \rC \in L^{\sm W}$ as an unordered bidirected
edge of $G^{\sm W}$ (under the \texttt{cdmg\_typed\_edges} encoding
announced in the statement block, $\lC w_{c_j}, w_{c_{j+1}} \rC$ is its
image under the quotient map $(V \sm W) \times (V \sm W) \twoheadrightarrow
\text{Sym}_2(V \sm W)$; the same comment applies to every
$\lC \cdot, \cdot \rC \in L^{\sm W}$ occurring below), and
\[
    a'_{k'-1} \;:=\; \text{the bidirected edge } \lC w'_{k'-1}, w'_{k'} \rC \in L^{\sm W}
\]
is a bidirected hinge edge for $p'$.

\emph{Sub-case (H.B): $p$'s hinge is directed with source $w_k \notin W$.}
Since $w_k \notin W$ and $c_j \le k - 1 < k$, the smallest $c$-index
strictly greater than $c_j$ is $k$ itself, i.e.\ $c_{j+1} = k$. By
clause~(e) for $p$ applied to the directed hinge, $1 \le k \le n - 1$;
with $c_s = n$ and the $(c_i)$ strictly increasing,
$k = c_{j+1} \le n - 1$ gives $j + 1 \le s - 1$, i.e.\
$k' = j + 1 \le s - 1 < s$. In particular the right arm of $p'$ is
non-trivial, so the boundary configuration $k' = s$ is incompatible with
sub-case~(H.B). Read the hinge-straddle sub-walk of $p$ \emph{in reverse}
from index $c_{j+1} = k$ down to index $c_j$: the traversed edges are
$p$'s hinge $a_{k - 1} = (w_k, w_{k - 1}) \in E$ (oriented from $w_k$ to
$w_{k - 1}$) followed by the left-arm-portion edges
$a_{k - 2}, a_{k - 3}, \dots, a_{c_j}$, each of the form
$a_i = (w_{i + 1}, w_i) \in E$ (oriented from $w_{i + 1}$ to $w_i$).
Concatenated in this order, they form a directed walk
\[
    \lp w_k,\, a_{k - 1},\, w_{k - 1},\, a_{k - 2},\, w_{k - 2},\, \dots,\, a_{c_j},\, w_{c_j} \rp
\]
in $G$ from $w_k = w'_{k'}$ to $w_{c_j} = w'_{k' - 1}$ of length
$k - c_j \ge 1$, with all intermediate vertices
$w_{k - 1}, \dots, w_{c_j + 1}$ in $W$ by maximality of the non-$W$
subsequence (every index strictly between $c_j$ and $c_{j+1} = k$ has
its vertex in $W$). By the contraction direction of $\Phi_E^{\sm W}$
(item~iii.\ of def \ref{def:G_marginalization}),
$(w_k, w_{c_j}) = (w'_{k'}, w'_{k' - 1}) \in E^{\sm W}$. Take
\[
    a'_{k'-1} \;:=\; (w'_{k'},\, w'_{k'-1}) \;\in\; E^{\sm W},
\]
a directed hinge edge for $p'$ realising clause~(d)'s directed
alternative at $k'$ with source $w'_{k'} = w_k$ at the interior
position. (This is structurally the same construction performed in
sub-claim~ii(b)'s $(\Rightarrow)$ direction, Region~H, reused here
because $p$'s hinge happens to have its source outside $W$ in this
sub-case.)

\emph{Verification of the bifurcation clauses for $p'$ in $G^{\sm W}$ with split index $k'$.}
\begin{itemize}
    \item \emph{Clause~(a).} $w'_0 = v_1 \ne v_2 = w'_s$ by clause~(a) for $p$.
    \item \emph{Clause~(b).} $w'_0 = v_1$ does not appear in $(w'_1, \dots, w'_s) \subseteq (w_{c_1}, \dots, w_{c_s}) \subseteq (w_1, \dots, w_n)$ by clause~(b) for $p$.
    \item \emph{Clause~(c).} Symmetrically, $w'_s = v_2$ does not appear in $(w'_0, \dots, w'_{s-1})$ by clause~(c) for $p$.
    \item \emph{Clause~(d).} Each $a'_j$ has the correct type by construction: $a'_j \in E^{\sm W}$ directed from $w'_{j+1}$ toward $w'_0$ for $j \le k' - 2$ (region L); $a'_j \in E^{\sm W}$ directed from $w'_j$ toward $w'_s$ for $j \ge k'$ (region R); and the hinge-straddle edge $a'_{k' - 1}$ is either the bidirected edge $\lC w'_{k' - 1}, w'_{k'} \rC \in L^{\sm W}$ (sub-case~(H.A)) or the directed edge $(w'_{k'}, w'_{k' - 1}) \in E^{\sm W}$ with source $w'_{k'}$ (sub-case~(H.B), including the refreshed Region~H produced when the pre-processing for degenerate identifications was applied), both of which are alternatives admitted by clause~(d) of def \ref{def:walks}, item~vi., applied to $p'$ at $k'$.
    \item \emph{Clause~(e).} The end-node arrowhead constraint at $w'_0 = v_1$: when $k' \ge 2$, $a'_0$ is the region~L directed edge $(w'_1, w'_0) \in E^{\sm W}$, automatically into $w'_0$; when $k' = 1$, $a'_0$ \emph{is} the hinge-straddle edge, and both admissible alternatives by clause~(d) --- the bidirected $\lC w'_0, w'_1 \rC \in L^{\sm W}$ (sub-case~(H.A)) and the directed $(w'_1, w'_0) \in E^{\sm W}$ (sub-case~(H.B)) --- are edges into $w'_0$, matching the clause~(e) case-analysis of def \ref{def:walks}, item~vi., at $k = 1$ (both alternatives at $k = 1$ are into $v_0$). The end-node arrowhead constraint at $w'_s = v_2$: when $k' \le s - 1$, $a'_{s - 1}$ is the region~R directed edge $(w'_{s - 1}, w'_s) \in E^{\sm W}$, automatically into $w'_s$ (sub-case~(H.B) forces $k' < s$, so this branch always applies in sub-case~(H.B)); when $k' = s$, $a'_{s - 1}$ \emph{is} the hinge-straddle edge, which forces sub-case~(H.A) (sub-case~(H.B) is incompatible with $k' = s$ as shown above) and the bidirected alternative $\lC w'_{s - 1}, w'_s \rC \in L^{\sm W}$, which is an edge into $w'_s$, matching the clause~(e) case-analysis of def \ref{def:walks}, item~vi., at $k = n$ (the bidirected alternative at $k = n$ is into $v_n$; the directed alternative at $k = n$ is excluded, which is consistent with sub-case~(H.B) being unreachable at $k' = s$).
\end{itemize}

Hence $p'$ together with split index $k'$ witnesses a bifurcation between $v_1$ and $v_2$ in $G^{\sm W}$, in the orientation $w'_0 = v_1$, $w'_s = v_2$. This is the LHS-to-RHS implication of sub-claim~ii(a).

\emph{Load-bearing corner case (the $n' = 1$ projected hinge).}
The construction above subsumes the corner case flagged in the rewritten statement's design notes. Consider the configuration where the original $p$ has $n = 2$, $k = 1$, hinge $a_0$ a directed alternative $(w_1, w_0) \in E$ with source $w_1$, and $w_1 \in W$ --- a $Y$-fork $v_1 \leftarrow w_1 \rightarrow v_2$ whose interior vertex is the sole interior vertex of $p$ and lies in $W$. The projection enumeration yields $c_0 = 0$, $c_1 = 2 = n$, so $s = 1$. The unique projected split index is $k' = 1$. The hinge straddle is the entire walk $p$, which is a bifurcation between $v_1$ and $v_2$ in $G$ with directed hinge and $w_1 \in W$ as the sole interior vertex; $\Phi_L^{\sm W}(v_1, v_2)$ holds via $p$ itself, and $v_1 \ne v_2$ by clause~(a) for $p$, so $\lC v_1, v_2 \rC \in L^{\sm W}$. The projected walk $p'$ collapses to the single bidirected edge $(v_1, \lC v_1, v_2 \rC, v_2)$ of length $n' = 1$ with $k' = 1$. \emph{This $n' = 1$ direct-bidirected-edge configuration of $p'$ is admitted as a valid bifurcation in $G^{\sm W}$ by def \ref{def:walks}, item~vi., precisely under the addition $[\text{\texttt{bifurcation\_index\_boundary\_excludes\_natural\_cases}}]$ to def \ref{def:walks}, item~vi., which authorises the $n = 1$, $k = 1$ boundary case of \emph{Walk.IsBifurcation} (a length-1 bidirected walk between distinct vertices).} Without that addition, the $(\Rightarrow)$ direction of sub-claim~ii(a) would fail on this configuration: the only candidate $G^{\sm W}$-witness for the RHS existential is the single bidirected edge $\lC v_1, v_2 \rC$, and excluding $n' = 1$ from \emph{Walk.IsBifurcation} would discard precisely that candidate.

\smallskip\noindent\emph{($\Leftarrow$) Converse direction (expansion of a bifurcation in $G^{\sm W}$).}
Let $p' = (w'_0, a'_0, \dots, a'_{s-1}, w'_s)$ together with split index $k' \in \{1, \dots, s\}$ witness a bifurcation between $v_1$ and $v_2$ in $G^{\sm W}$ per def \ref{def:walks}, item~vi., applied to $G^{\sm W}$. So $s \ge 1$, $\lC w'_0, w'_s \rC = \lC v_1, v_2 \rC$, and clauses~(a)--(e) hold for $p'$ and $k'$, with edge-membership read against $E^{\sm W}$ and $L^{\sm W}$. WLOG $w'_0 = v_1$, $w'_s = v_2$.

\emph{Expand each non-hinge arm edge.}
For each $j \in \{0, 1, \dots, s-1\} \sm \lC k' - 1 \rC$, the edge $a'_j$ lies in $E^{\sm W}$ (by clause~(d) for $p'$, in the left-arm or right-arm orientation).
\begin{itemize}
    \item If $j \le k' - 2$ (left arm of $p'$): $a'_j = (w'_{j+1}, w'_j) \in E^{\sm W}$. By the expansion direction of $\Phi_E^{\sm W}$, there exists a directed walk $q_j$ in $G$ from $w'_{j+1}$ to $w'_j$ of length $\ge 1$ with all intermediate vertices in $W$.
    \item If $j \ge k'$ (right arm of $p'$): $a'_j = (w'_j, w'_{j+1}) \in E^{\sm W}$. By the expansion direction of $\Phi_E^{\sm W}$, there exists a directed walk $q_j$ in $G$ from $w'_j$ to $w'_{j+1}$ of length $\ge 1$ with intermediate vertices in $W$.
\end{itemize}

\emph{Expand the hinge edge.}
Two sub-cases by the type of $a'_{k'-1}$ (clause~(d) for $p'$ admits both):
\begin{itemize}
    \item \emph{Bidirected hinge of $p'$:} $\lC w'_{k'-1}, w'_{k'} \rC \in L^{\sm W}$. By the expansion direction of $\Phi_L^{\sm W}$ (item~iv.\ of def \ref{def:G_marginalization}), there exists a bifurcation $q_{k'-1}$ in $G$ between $w'_{k'-1}$ and $w'_{k'}$ with all intermediate vertices in $W$ --- explicitly, a walk together with a split index satisfying clauses~(a)--(e) of def \ref{def:walks}, item~vi.\ (the LN's converse-direction reading: ``a bifurcation in $G$ of the form $w'_{k'-1} \hut z_1 \hut \dots \hut z_{\tilde k - 1} \hus z_{\tilde k} \tuh \dots \tuh z_{\tilde n - 1} \tuh w'_{k'}$ with $z_1, \dots, z_{\tilde n - 1} \in W$''). (The bidirected hinge $\lC w'_{k'-1}, w'_{k'} \rC$ is well-defined as a $\text{Sym}_2$ element under the \texttt{cdmg\_typed\_edges} encoding announced in the statement block; the expansion existential is consequently invariant under swapping $w'_{k'-1}$ and $w'_{k'}$, matching the symmetric ``bifurcation between $\ul{v}$ and $\ol{v}$'' phrasing of $\Phi_L^{\sm W}$.)
    \item \emph{Directed hinge of $p'$:} $a'_{k'-1} = (w'_{k'}, w'_{k'-1}) \in E^{\sm W}$ (the source-of-$p'$ case). By the expansion direction of $\Phi_E^{\sm W}$, there exists a directed walk $q_{k'-1}$ in $G$ from $w'_{k'}$ to $w'_{k'-1}$ of length $\ge 1$ with intermediate vertices in $W$.
\end{itemize}

\emph{Concatenate.}
Concatenate $q_0, q_1, \dots, q_{s-1}$ along their shared endpoints $w'_1, w'_2, \dots, w'_{s-1}$ in $G$. The resulting walk $q$ in $G$ has the bifurcation shape: its left arm (the concatenation of $q_0, q_1, \dots, q_{k'-2}$, each a directed walk in $G$ ending at $w'_j$ for $j = 0, 1, \dots, k'-1$ in sequence) reads as a chain of directed edges in $E$ flowing from interior toward $w'_0 = v_1$; its hinge segment ($q_{k'-1}$) is either a bifurcation in $G$ between $w'_{k'-1}$ and $w'_{k'}$ (bidirected sub-case) or a directed walk from $w'_{k'}$ to $w'_{k'-1}$ (directed sub-case, with $w'_{k'}$ as source); and its right arm (the concatenation of $q_{k'}, \dots, q_{s-1}$) reads as a chain of directed edges flowing toward $w'_s = v_2$.

\emph{Verification of the bifurcation clauses for $q$ in $G$ with the induced hinge index.}
Let $K$ denote the position in $q$'s vertex list of the hinge of $q_{k'-1}$ (in the bidirected sub-case, the hinge position within the bifurcation $q_{k'-1}$; in the directed sub-case, the position of $w'_{k'}$). Clauses~(a)--(e) of def \ref{def:walks}, item~vi., applied to $q$ and $K$:
\begin{itemize}
    \item Clause~(a): $q$'s endpoints are $w'_0 = v_1$ and $w'_s = v_2$; $v_1 \ne v_2$ by clause~(a) for $p'$.
    \item Clauses~(b), (c): the endpoints $v_1, v_2$ of $q$ do not repeat as interior vertices because the interior vertices of $q$ consist of (i) the interior vertices of the expansions $q_j$, which all lie in $W$ (so cannot equal $v_1, v_2 \in (J \cup V) \sm W$), and (ii) the projected vertices $w'_1, w'_2, \dots, w'_{s-1}$ of $p'$, none of which equals $w'_0 = v_1$ or $w'_s = v_2$ by clauses~(b) and~(c) for $p'$.
    \item Clauses~(d), (e): the directed structure of the left- and right-arm expansions $q_j$ ($j \ne k' - 1$) provides the correct arrowhead orientation toward $v_1$ on the left arm and toward $v_2$ on the right arm; the hinge segment $q_{k'-1}$ provides the bifurcation hinge (either bidirected from the bidirected sub-case, or directed-with-source-$w'_{k'}$ from the directed sub-case).
\end{itemize}

Hence $q$ together with the induced hinge index witnesses a bifurcation between $v_1$ and $v_2$ in $G$ per def \ref{def:walks}, item~vi. This is the RHS-to-LHS implication of sub-claim~ii(a).

\medskip\noindent\textbf{Sub-claim~ii(b): sourced bifurcation biconditional.}

Let $v_1, v_2, v_3 \in (J \cup V) \sm W = J \cup (V \sm W)$. We follow the same projection / expansion strategy as for sub-claim~ii(a), additionally tracking the source $v_3$ at both ends of the biconditional.

\smallskip\noindent\emph{($\Rightarrow$) Forward direction (contraction, with source-tracking).}
Let $p$, $k$ witness a bifurcation between $v_1$ and $v_2$ with source $v_3$ in $G$: $p$ and $k$ satisfy clauses~(a)--(e) of def \ref{def:walks}, item~vi., and clause~(d) is realised at $k$ by the directed alternative $a_{k-1} = (w_k, w_{k-1}) \in E$ with $w_k = v_3$. By clause~(e), $1 \le k \le n - 1$, so $n \ge 2$ and $v_3 = w_k$ is an interior vertex of $p$ distinct from $v_1, v_2$.

WLOG $w_0 = v_1$, $w_n = v_2$ (the other orientation is handled symmetrically).

\emph{The source survives the projection.} By hypothesis $v_3 \in (J \cup V) \sm W$, so $w_k = v_3 \notin W$. In the maximal-non-$W$-subsequence enumeration $0 = c_0 < c_1 < \dots < c_s = n$ from sub-claim~ii(a)'s $(\Rightarrow)$ direction, the index $k$ appears as some $c_{j_0}$ with $w'_{j_0} = w_{c_{j_0}} = w_k = v_3$. Since $w_k$ is an interior vertex distinct from $v_1, v_2$ (clauses~(a), (b), (c), (e) for $p$), we have $1 \le j_0 \le s - 1$.

\emph{The projected split index coincides with the source position.} The unique projected split index $k'$ from sub-claim~ii(a)'s $(\Rightarrow)$ direction satisfies $c_{k'-1} \le k - 1$ and $c_{k'} \ge k$. Since $c_{j_0} = k$, we have $c_{j_0 - 1} < c_{j_0} = k$ (so $c_{j_0 - 1} \le k - 1$) and $c_{j_0} = k \ge k$, so by uniqueness $k' = j_0$. In particular, $w'_{k'} = w'_{j_0} = v_3$.

\emph{The projected hinge is a directed edge with source $v_3$.}
The hinge-straddling sub-case (region~H) from sub-claim~ii(a)'s $(\Rightarrow)$ direction yielded a bidirected hinge $\lC w'_{k'-1}, w'_{k'} \rC \in L^{\sm W}$ generically. For sub-claim~ii(b) we strengthen this to a \emph{directed} hinge in $E^{\sm W}$: read the hinge-straddling sub-walk of $p$ from index $c_{k'-1}$ to index $c_{k'} = k$ \emph{in the reverse order} (from index $k$ down to $c_{k'-1}$). This reversed traversal consists of, in order: the hinge edge $a_{k-1} = (w_k, w_{k-1}) \in E$ (directed from $w_k = v_3$ to $w_{k-1}$), followed by the left-arm edges $a_{k-2}, a_{k-3}, \dots, a_{c_{k'-1}}$ each of the form $a_i = (w_{i+1}, w_i) \in E$ (directed from $w_{i+1}$ to $w_i$). Concatenated in this order, these are the edges of a directed walk
\[
    \lp w_k,\, a_{k-1},\, w_{k-1},\, a_{k-2},\, w_{k-2},\, \dots,\, a_{c_{k'-1}},\, w_{c_{k'-1}} \rp
\]
in $G$ from $w_k = v_3$ to $w_{c_{k'-1}} = w'_{k'-1}$ of length $k - c_{k'-1} \ge 1$, with all intermediate vertices $w_{k-1}, w_{k-2}, \dots, w_{c_{k'-1} + 1}$ in $W$ by maximality of the non-$W$ subsequence. By the contraction direction of $\Phi_E^{\sm W}$, $(w_k, w_{c_{k'-1}}) = (v_3, w'_{k'-1}) \in E^{\sm W}$. Take
\[
    a'_{k'-1} \;:=\; (w'_{k'},\, w'_{k'-1}) \;=\; (v_3, w'_{k'-1}) \;\in\; E^{\sm W},
\]
a directed hinge edge of $p'$ with $w'_{k'} = v_3$ in the source-of-hinge interior position.

\emph{The non-hinge edges.} For $j \le k' - 2$ (left arm) and $j \ge k'$ (right arm), construct $a'_j$ exactly as in sub-claim~ii(a)'s $(\Rightarrow)$ direction (regions L and R, both yielding $E^{\sm W}$-edges).

\emph{Verification.} Clauses~(a)--(e) of def \ref{def:walks}, item~vi., for $p'$ and $k'$ applied to $G^{\sm W}$ are verified exactly as in sub-claim~ii(a)'s $(\Rightarrow)$ direction's verification, with the strengthening that clause~(d) at the hinge $k'$ is realised by the \emph{directed} alternative $a'_{k'-1} = (w'_{k'}, w'_{k'-1}) \in E^{\sm W}$ with $w'_{k'} = v_3$. Hence $p'$ and $k'$ witness a bifurcation between $v_1$ and $v_2$ with source $v_3$ in $G^{\sm W}$ per def \ref{def:walks}, item~vi., and its closing paragraph.

\smallskip\noindent\emph{($\Leftarrow$) Converse direction (expansion, with source-tracking).}
Let $p' = (w'_0, a'_0, \dots, a'_{s-1}, w'_s)$ together with split index $k' \in \{1, \dots, s\}$ witness a bifurcation between $v_1$ and $v_2$ with source $v_3$ in $G^{\sm W}$. So clauses~(a)--(e) hold for $p'$ and $k'$ applied to $G^{\sm W}$, and clause~(d) at $k'$ is realised by the directed alternative $a'_{k'-1} = (w'_{k'}, w'_{k'-1}) \in E^{\sm W}$ with $w'_{k'} = v_3$. By clause~(e) for $p'$, $1 \le k' \le s - 1$, so $s \ge 2$ and $v_3 = w'_{k'}$ is an interior vertex of $p'$ distinct from $v_1, v_2$.

WLOG $w'_0 = v_1$, $w'_s = v_2$.

\emph{Expand non-hinge edges.} Construct directed walks $q_j$ in $G$ for $j \in \{0, 1, \dots, s-1\} \sm \lC k' - 1 \rC$ exactly as in sub-claim~ii(a)'s $(\Leftarrow)$ direction.

\emph{Expand the hinge edge.} Since $a'_{k'-1} = (w'_{k'}, w'_{k'-1}) = (v_3, w'_{k'-1}) \in E^{\sm W}$, by the expansion direction of $\Phi_E^{\sm W}$ there exists a directed walk $q_{k'-1}$ in $G$ from $v_3 = w'_{k'}$ to $w'_{k'-1}$ of length $\ge 1$ with all intermediate vertices in $W$.

\emph{Concatenate.} Concatenate $q_0, q_1, \dots, q_{s-1}$ along their shared endpoints to obtain a walk $q$ in $G$ from $v_1$ to $v_2$. Its left arm (concatenation of $q_0, q_1, \dots, q_{k'-1}$) ends at the vertex sequence reaching $v_3 = w'_{k'}$ via $q_{k'-1}$; reading the directions of $q_{k'-1}$ (a directed walk \emph{starting} at $v_3$), the leftmost edge of $q_{k'-1}$ in $q$'s vertex order is the directed edge $(v_3, z_1) \in E$ where $z_1$ is the second vertex of $q_{k'-1}$ --- equivalently, the hinge of $q$ is at the position of $v_3$ in $q$'s vertex list, realised by the directed alternative whose interior position is $v_3$.

\emph{Verification.} Clauses~(a)--(e) of def \ref{def:walks}, item~vi., for $q$ and its induced hinge index applied to $G$ are verified as in sub-claim~ii(a)'s $(\Leftarrow)$ direction's verification, with the strengthening that clause~(d) at the hinge of $q$ is realised by the directed alternative whose interior vertex is $v_3$. Hence $q$ and its hinge index witness a bifurcation between $v_1$ and $v_2$ with source $v_3$ in $G$ per def \ref{def:walks}, item~vi.

This completes sub-claim~ii(b).

\medskip\noindent\textbf{Sub-claim~iii(a): preservation of acyclicity.}

Suppose $G$ is acyclic in the sense of def \ref{def-acylic}, i.e.\ for every $v \in J \cup V$, there does not exist a non-trivial directed walk from $v$ to $v$ in $G$ (where ``non-trivial'' means walk length $n \ge 1$). We show $G^{\sm W}$ is acyclic in the sense of def \ref{def-acylic}, evaluated on the carrier $J \cup (V \sm W)$.

Suppose for contradiction that some $v \in J \cup (V \sm W)$ and some directed walk $p' = (u'_0, a'_0, u'_1, \dots, a'_{s-1}, u'_s)$ in $G^{\sm W}$ with $u'_0 = u'_s = v$ have length $s \ge 1$ --- a non-trivial directed cycle in $G^{\sm W}$.

By the expansion direction of $\Phi_E^{\sm W}$ (the same lifting step as in sub-claim~i's $(\Leftarrow)$ direction): for each $j \in \{0, 1, \dots, s-1\}$, $a'_j = (u'_j, u'_{j+1}) \in E^{\sm W}$ admits a directed walk $q_j$ in $G$ from $u'_j$ to $u'_{j+1}$ of length $|q_j| \ge 1$ with all intermediate vertices in $W$. (The length-$\ge 1$ bound is the third conjunct of $\Phi_E^{\sm W}$, valid for every $E^{\sm W}$-edge regardless of whether $u'_j = u'_{j+1}$; see also the witnessing-walk discussion of the setup paragraph.)

Concatenate $q_0, q_1, \dots, q_{s-1}$ along their shared endpoints $u'_1, u'_2, \dots, u'_{s-1}$ in $G$ to obtain a directed walk $q$ in $G$ from $u'_0 = v$ to $u'_s = v$ of length
\[
    |q| \;=\; \sum_{j=0}^{s-1} |q_j| \;\ge\; \sum_{j=0}^{s-1} 1 \;=\; s \;\ge\; 1.
\]
Hence $q$ is a non-trivial directed walk from $v$ to itself in $G$, with $v \in J \cup (V \sm W) \subseteq J \cup V$. This contradicts def \ref{def-acylic}, which forbids any non-trivial directed walk from a vertex to itself in $G$.

Hence no non-trivial directed cycle in $G^{\sm W}$ exists, so $G^{\sm W}$ is acyclic in the sense of def \ref{def-acylic}. This establishes sub-claim~iii(a).

\medskip\noindent\textbf{Sub-claim~iii(b): topological order is induced by restriction.}

Let $<_G$ be a binary relation on $J \cup V$ that is a topological order of $G$ in the sense of def \ref{def-topological-order}. Explicitly:
\begin{itemize}
    \item $<_G$ is a strict total order on $J \cup V$ (irreflexive, transitive, and trichotomous on $J \cup V$);
    \item For every $v_1, v_2 \in J \cup V$, $v_1 \in \Pa^G(v_2) \Rightarrow v_1 <_G v_2$ (parent precedence, where $\Pa^G(v_2)$ is the parent set in $G$ per def \ref{def_3_5}, item~i.).
\end{itemize}
Let $<_{G^{\sm W}}$ be the restriction of $<_G$ to the smaller carrier $J \cup (V \sm W)$ as defined in the statement block above: for every $v_1, v_2 \in J \cup (V \sm W)$, $v_1 <_{G^{\sm W}} v_2 :\Longleftrightarrow v_1 <_G v_2$. We verify the four conjuncts of def \ref{def-topological-order} applied to $G^{\sm W}$ (strict-total-order conjuncts plus parent precedence) for $<_{G^{\sm W}}$.

\smallskip
We use throughout the carrier inclusion
\[
    J \cup (V \sm W) \;\subseteq\; J \cup V,
\]
which holds since $V \sm W \subseteq V$.

\emph{Irreflexivity of $<_{G^{\sm W}}$.}
Let $v \in J \cup (V \sm W) \subseteq J \cup V$. Irreflexivity of $<_G$ on $J \cup V$ at $v$ gives $\neg (v <_G v)$, i.e.\ $\neg (v <_{G^{\sm W}} v)$ by the pointwise definition.

\emph{Transitivity of $<_{G^{\sm W}}$.}
Let $v_1, v_2, v_3 \in J \cup (V \sm W) \subseteq J \cup V$ with $v_1 <_{G^{\sm W}} v_2$ and $v_2 <_{G^{\sm W}} v_3$. The pointwise definition gives $v_1 <_G v_2$ and $v_2 <_G v_3$; transitivity of $<_G$ on $J \cup V$ at $v_1, v_2, v_3$ gives $v_1 <_G v_3$, i.e.\ $v_1 <_{G^{\sm W}} v_3$ by the pointwise definition.

\emph{Trichotomy of $<_{G^{\sm W}}$.}
Let $v_1, v_2 \in J \cup (V \sm W) \subseteq J \cup V$. Trichotomy of $<_G$ on $J \cup V$ at $v_1, v_2$ gives exactly one of $v_1 <_G v_2$, $v_1 = v_2$, $v_2 <_G v_1$. By the pointwise definition, these read as exactly one of $v_1 <_{G^{\sm W}} v_2$, $v_1 = v_2$, $v_2 <_{G^{\sm W}} v_1$.

\smallskip
Combining the three properties, $<_{G^{\sm W}}$ is a strict total order on $J \cup (V \sm W)$.

\emph{Parent precedence of $<_{G^{\sm W}}$ with respect to $G^{\sm W}$.}
Let $v_1, v_2 \in J \cup (V \sm W)$ with $v_1 \in \Pa^{G^{\sm W}}(v_2)$. By def \ref{def_3_5}, item~i., applied to $G^{\sm W}$, this is $(v_1, v_2) \in E^{\sm W}$. By the expansion direction of $\Phi_E^{\sm W}$ (item~iii.\ of def \ref{def:G_marginalization}), there exists a directed walk
\[
    p \;=\; (u_0, a_0, u_1, a_1, \dots, a_{n-1}, u_n)
\]
in $G$ from $v_1 = u_0$ to $v_2 = u_n$ of length $n \ge 1$ with all intermediate vertices $u_1, \dots, u_{n-1} \in W$.

For each step $i \in \{0, 1, \dots, n-1\}$, $a_i = (u_i, u_{i+1}) \in E$, so $u_i \in \Pa^G(u_{i+1})$ by def \ref{def_3_5}, item~i., applied to $G$. The vertices $u_0, u_1, \dots, u_n$ are all in $J \cup V$ (walk vertices of $G$ are in the carrier of $G$, per def \ref{def-cdmg}). By parent precedence of $<_G$ on $G$ applied at each step, $u_i <_G u_{i+1}$ for every $i \in \{0, 1, \dots, n-1\}$.

We thus have a chain
\[
    u_0 \,<_G\, u_1 \,<_G\, u_2 \,<_G\, \cdots \,<_G\, u_n
\]
in $J \cup V$ with $n \ge 1$ links. By iterated application of transitivity of $<_G$ on $J \cup V$ along this chain (immediate from $u_0 <_G u_1$ when $n = 1$; by induction on $n$ otherwise), $u_0 <_G u_n$, i.e.\ $v_1 <_G v_2$.

Since $v_1, v_2 \in J \cup (V \sm W)$, the pointwise definition of $<_{G^{\sm W}}$ gives $v_1 <_{G^{\sm W}} v_2$. This is the parent-precedence clause of $<_{G^{\sm W}}$ on $G^{\sm W}$.

\smallskip
Combining the four conjuncts (irreflexivity, transitivity, trichotomy, parent precedence), $<_{G^{\sm W}}$ is a topological order of $G^{\sm W}$ in the sense of def \ref{def-topological-order}. This establishes sub-claim~iii(b) and completes the proof of the remark.
\end{proof}

\end{document}
